\answerpaper{数学 A 类}

\answerproblem{一、}{15}
记三维欧氏空间中不共面的四点 $A,B,C,D$ 生成的四面体为 $ABCD$．在四面体 $ABCD$ 内部取一点 $O$，设四面体 $OBCD$ 的体积为 $V_A$，四面体 $OACD$ 的体积为 $V_B$，四面体 $OABD$ 的体积为 $V_C$，四面体 $OABC$ 的体积为 $V_D$．（图略：四面体 $ABCD$，内部一点 $O$ 与四个顶点分别用虚线相连）
\begin{enumerate}[label=(\arabic*)]
  \item 求证：$\Bigl(\overrightarrow{OA}\cdot\bigl(\overrightarrow{OC}\times\overrightarrow{OD}\bigr)\Bigr)\Bigl(\overrightarrow{OB}\cdot\bigl(\overrightarrow{OC}\times\overrightarrow{OD}\bigr)\Bigr)<0$；
  \item 求证：$V_A\overrightarrow{OA}+V_B\overrightarrow{OB}+V_C\overrightarrow{OC}+V_D\overrightarrow{OD}=\vec{0}$．
\end{enumerate}
\begin{solution}
\textbf{(1)} 因为点 $A,B$ 落在 $O,C,D$ 所决定的平面的两侧，若 $\overrightarrow{OA},\overrightarrow{OC},\overrightarrow{OD}$ 为右手系，则 $\overrightarrow{OB},\overrightarrow{OC},\overrightarrow{OD}$ 为左手系，反之亦然．故
\[
\Bigl(\overrightarrow{OA}\cdot\bigl(\overrightarrow{OC}\times\overrightarrow{OD}\bigr)\Bigr)\Bigl(\overrightarrow{OB}\cdot\bigl(\overrightarrow{OC}\times\overrightarrow{OD}\bigr)\Bigr)<0.
\]

\textbf{(2)} 三维欧氏空间中四个向量 $\overrightarrow{OA},\overrightarrow{OB},\overrightarrow{OC},\overrightarrow{OD}$ 线性相关．因为 $O$ 点在四面体 $ABCD$ 内部，所以这四个向量中每三个向量线性无关，故存在非零实数 $x,y,z,w$ 使得
\begin{equation*}
x\overrightarrow{OA}+y\overrightarrow{OB}+z\overrightarrow{OC}+w\overrightarrow{OD}=\vec{0}.\tag{$*$}
\end{equation*}
将 $(*)$ 式与向量 $\overrightarrow{OC}\times\overrightarrow{OD}$ 作内积，得到
\[
x\,\overrightarrow{OA}\cdot\bigl(\overrightarrow{OC}\times\overrightarrow{OD}\bigr)
+y\,\overrightarrow{OB}\cdot\bigl(\overrightarrow{OC}\times\overrightarrow{OD}\bigr)=0.
\]
不妨设 $\overrightarrow{OA},\overrightarrow{OC},\overrightarrow{OD}$ 为右手系，则 $\overrightarrow{OB},\overrightarrow{OC},\overrightarrow{OD}$ 为左手系，而三个共起点向量决定的平行六面体的体积等于 $6$ 倍的这三个向量起点和三终点顶点决定的四面体体积，故
\[
\overrightarrow{OA}\cdot\bigl(\overrightarrow{OC}\times\overrightarrow{OD}\bigr)=6V_B,\qquad
\overrightarrow{OB}\cdot\bigl(\overrightarrow{OC}\times\overrightarrow{OD}\bigr)=-6V_A,\qquad
y=\frac{V_B}{V_A}x.
\]
同理，将 $(*)$ 式与向量 $\overrightarrow{OB}\times\overrightarrow{OD}$ 作内积，得到 $z=\dfrac{V_C}{V_A}x$；将 $(*)$ 式与向量 $\overrightarrow{OB}\times\overrightarrow{OC}$ 作内积，得到 $w=\dfrac{V_D}{V_A}x$．代入 $(*)$ 式，得到
\[
V_A\overrightarrow{OA}+V_B\overrightarrow{OB}+V_C\overrightarrow{OC}+V_D\overrightarrow{OD}=\vec{0}.
\]
\end{solution}

\answerproblem{二、}{15}
设 $N\geqslant 1$，$\mathcal{S}$ 是包含 $N$ 个整数的集合，满足如下的加性唯一性条件：

若 $n_k\in\mathcal{S}\ (1\leqslant k\leqslant 4)$ 且 $n_1+n_2=n_3+n_4$，则必有 $n_1=n_3$ 或 $n_1=n_4$．

令 $f(x):=\sum\limits_{n\in\mathcal{S}}\ee^{2\pi\ii nx}$．
\begin{enumerate}[label=(\arabic*)]
  \item 计算 $\displaystyle\int_0^1|f(x)|^2\dd x$，$\displaystyle\int_0^1|f(x)|^4\dd x$．
  \item 证明：$\displaystyle\int_0^1|f(x)|\dd x\geqslant\frac{1}{2\sqrt{N}}$．
\end{enumerate}
\begin{solution}
\textbf{(1)} 我们有
\[
|f(x)|^2=f(x)\overline{f(x)}
=\Bigl(\sum_{m\in\mathcal{S}}\ee^{2\pi\ii mx}\Bigr)\Bigl(\sum_{n\in\mathcal{S}}\ee^{-2\pi\ii nx}\Bigr)
\equiv N+g(x),
\]
其中
\[
g(x)=\sum_{\substack{m,n\in\mathcal{S}\\ m\neq n}}\ee^{2\pi\ii(m-n)x}.
\]
于是，注意到当整数 $k\neq 0$ 时，$\displaystyle\int_0^1\ee^{2\pi\ii kx}\dd x=0$，立即得到 $\displaystyle\int_0^1 g(x)\dd x=0$，从而
\[
\int_0^1|f(x)|^2\dd x=N.
\]
而由题设条件，对于 $m_1\neq n_1$，$m_2\neq n_2$，若 $m_1\neq m_2$，则 $m_1-n_1\neq m_2-n_2$．于是，
\[
g(x)=\sum_{k\in\mathcal{S}_g}\ee^{2\pi\ii kx},
\]
这里集合 $\mathcal{S}_g$ 由 $N^2-N$ 个两两不同的非零整数组成．这样，又有
\[
\int_0^1|f(x)|^4\dd x
=\int_0^1\bigl(N^2+2Ng(x)+|g(x)|^2\bigr)\dd x
=N^2+0+(N^2-N)=2N^2-N.
\]

\textbf{(2)} 对于本小题，我们提供两种方式的证明．

\textbf{法 I.} 注意到 $\displaystyle\int_0^1 g(x)\dd x=0$，我们有
\begin{align*}
\sqrt{N}-\int_0^1|f(x)|\dd x
&=\int_0^1\bigl(\sqrt{N}-\sqrt{N+g(x)}\bigr)\dd x
=-\int_0^1\frac{g(x)}{\sqrt{N}+\sqrt{N+g(x)}}\dd x\\
&=\int_0^1\Bigl(\frac{g(x)}{2\sqrt{N}}-\frac{g(x)}{\sqrt{N}+\sqrt{N+g(x)}}\Bigr)\dd x\\
&=\int_0^1\frac{g^2(x)}{2\sqrt{N}\bigl(\sqrt{N}+\sqrt{N+g(x)}\bigr)^2}\dd x\\
&\leqslant\int_0^1\frac{g^2(x)}{2N\sqrt{N}}\dd x
=\frac{N^2-N}{2N\sqrt{N}}<\frac{\sqrt{N}}{2}.
\end{align*}
因此，
\[
\int_0^1|f(x)|\dd x>\frac{\sqrt{N}}{2}.
\]

\textbf{法 II.} 由 H\"older 不等式，我们有
\begin{align*}
N=\int_0^1|f(x)|^2\dd x
&\leqslant\Bigl(\int_0^1|f(x)|\dd x\Bigr)^{\frac23}\Bigl(\int_0^1|f(x)|^4\dd x\Bigr)^{\frac13}\\
&=\Bigl(\int_0^1|f(x)|\dd x\Bigr)^{\frac23}\bigl(2N^2-N\bigr)^{\frac13}.
\end{align*}
因此，
\[
\int_0^1|f(x)|\dd x\geqslant\sqrt{\frac{N^3}{2N^2-N}}>\frac{\sqrt{N}}{\sqrt{2}}>\frac{\sqrt{N}}{2}.
\]
\end{solution}

\answerproblem{三、}{20}
设 $m,n$ 为大于 $2$ 的整数，$a_1,\cdots,a_{m+1}$ 为任意 $m+1$ 个有理数，$\mathbb{Q}^{n\times n}$ 为有理数域上 $n$ 阶方阵全体．证明：
\begin{enumerate}[label=(\arabic*)]
  \item $\mathbb{Q}^{n\times n}$ 中存在 $m$ 个元素 $B_1,\cdots,B_m$ 使得行列式 $|B_j|=j\ (j=1,\cdots,m)$ 成立．
  \item $\mathbb{Q}^{n\times n}$ 中存在 $m$ 个元素 $A_1,\cdots,A_m$ 使得下列两条同时成立：
  \begin{enumerate}[label=(\roman*)]
    \item $|A_j|=a_j\ (j=1,\cdots,m)$；
    \item $|A_1-A_2-\cdots-A_m|=a_{m+1}$．
  \end{enumerate}
\end{enumerate}
\begin{solution}
\textbf{(1)} 令
\[
B_1=\begin{pmatrix}
1& & &\\
& \ddots& &\\
& & \ddots&\\
& & &1
\end{pmatrix},\quad
B_2=\begin{pmatrix}
1& & &\\
&2& &\\
& & \ddots&\\
& & &1
\end{pmatrix},\quad\cdots,\quad
B_m=\begin{pmatrix}
1& & &\\
& \ddots& &\\
& &1&\\
& & &m
\end{pmatrix},
\]
则有
\[
|B_j|=j,\qquad \forall\, j=1,\cdots,m.
\]

\textbf{(2)} 令
\[
A_1=\begin{pmatrix}
a_1&0& &\\
1&1& &\\
& & \ddots&\\
& & &1
\end{pmatrix},\quad
A_2=\begin{pmatrix}
a_2&b& &\\
0&1& &\\
& & \ddots&\\
& & &1
\end{pmatrix},
\]
\[
A_i=\begin{pmatrix}
a_i& & &\\
&1& &\\
& & \ddots&\\
& & &1
\end{pmatrix},\qquad i=3,\cdots,m.
\]
则有
\begin{enumerate}[label=(\roman*)]
  \item $|A_j|=a_j$，$j=1,\cdots,m$；
  \item
  \[
  |A_1-A_2-\cdots-A_m|=
  \begin{vmatrix}
  a_1-\sum\limits_{i=2}^{m}a_i&-b& &\\
  1&-(m-2)& &\\
  & & \ddots&\\
  & & &-(m-2)
  \end{vmatrix}.
  \]
\end{enumerate}
于是有
\[
|A_1-A_2-\cdots-A_m|=\bigl(-s(m-2)+b\bigr)\bigl[-(m-2)\bigr]^{m-2},
\]
其中
\[
s=a_1-\sum_{i=2}^{m}a_i.
\]
取
\[
b=s(m-2)+\frac{(-1)^m a_{m+1}}{(m-2)^{m-2}},
\]
则有 $b\in\mathbb{Q}$，从而 $A_1,\cdots,A_m$ 满足条件．证毕．
\end{solution}

\answerproblem{四、}{15}
设 $K$ 为数域，$\mathcal{A}:V\to V$ 是 $n$ 维 $K$-向量空间 $V$ 上的线性算子（线性变换），$\mu_{\mathcal{A}}(x)$，$\chi_{\mathcal{A}}(x)\in K[x]$ 分别是 $\mathcal{A}$ 的极小多项式和特征多项式．证明：
\begin{enumerate}[label=(\arabic*)]
  \item 存在 $\alpha\in V$ 使 $\mu_{\mathcal{A}}(x)=\mu_{\mathcal{A},\alpha}(x)$，其中 $\mu_{\mathcal{A},\alpha}(x)$ 是集合
  \[
    \{\,f(x)\in K[x]\mid f(\mathcal{A})(\alpha)=0\,\}
  \]
  中首项系数为 1，次数最小的多项式；
  \item $K[\mathcal{A}]\cdot\alpha:=\{\,f(\mathcal{A})(\alpha)\mid \forall f(x)\in K[x]\,\}$ 是 $\deg\mu_{\mathcal{A},\alpha}(x)$ 维 $\mathcal{A}$-不变子空间；
  \item 设 $\mathcal{B}:V\to V$ 是任意线性算子．若 $\mu_{\mathcal{A}}(x)=\chi_{\mathcal{A}}(x)$，则
  \[
    \mathcal{A}\cdot\mathcal{B}=\mathcal{B}\cdot\mathcal{A}
    \iff \text{存在}\ f(x)\in K[x]\ \text{使}\ \mathcal{B}=f(\mathcal{A}).
  \]
\end{enumerate}
\begin{solution}
\textbf{(1)} 设 $\mu_{\mathcal{A}}(x)=p_1(x)^{m_1}\cdots p_s(x)^{m_s}$，$m_i>0$，是 $\mu_{\mathcal{A}}(x)$ 的不可约分解，则 $V=V_1\oplus\cdots\oplus V_s$，$\mathcal{A}(V_i)\subset V_i$，且 $\mathcal{A}_i=\mathcal{A}|_{V_i}$ 的极小多项式 $\mu_{\mathcal{A}_i}(x)=p_i(x)^{m_i}$．先证明存在 $\alpha_i\in V_i$ 使 $\mu_{\mathcal{A},\alpha_i}(x)=p_i(x)^{m_i}$．事实上，若 $m_i=1$，$\forall\alpha_i\in V_i$ 非零，则 $\mu_{\mathcal{A},\alpha_i}(x)\mid\mu_{\mathcal{A}_i}(x)=p_i(x)$，故 $\mu_{\mathcal{A},\alpha_i}(x)=p_i(x)$．若 $m_i>1$，则 $p_i(\mathcal{A})^{m_i-1}\neq 0$，因此存在 $\alpha_i\in V_i$ 使 $p_i(\mathcal{A})^{m_i-1}\alpha_i\neq 0$．另一方面，由 $\mu_{\mathcal{A},\alpha_i}(x)\mid\mu_{\mathcal{A}_i}(x)$ 得 $\mu_{\mathcal{A},\alpha_i}(x)=p_i(x)^{m_i'}$（$m_i'\leqslant m_i$）．若 $m_i'<m_i$，则
\[
p_i(\mathcal{A})^{m_i-1}\alpha_i=p_i(\mathcal{A})^{m_i-m_i'-1}p_i(\mathcal{A})^{m_i'}\alpha_i=0,
\]
此与 $p_i(\mathcal{A})^{m_i-1}\alpha_i\neq 0$ 矛盾，故 $\mu_{\mathcal{A},\alpha_i}(x)=p_i(x)^{m_i}$．令 $\alpha=\alpha_1+\cdots+\alpha_s$，则
\[
0=\mu_{\mathcal{A},\alpha}(\mathcal{A})\alpha
=\mu_{\mathcal{A},\alpha}(\mathcal{A})\alpha_1+\cdots+\mu_{\mathcal{A},\alpha}(\mathcal{A})\alpha_s,
\qquad \mu_{\mathcal{A},\alpha}(\mathcal{A})\alpha_i\in V_i.
\]
由 $V_1+\cdots+V_s$ 是直和得
\[
\mu_{\mathcal{A},\alpha}(\mathcal{A})\alpha_i=0,\qquad 1\leqslant i\leqslant s,
\]
从而 $p_i(x)^{m_i}\mid\mu_{\mathcal{A},\alpha}(x)$，$1\leqslant i\leqslant s$，由此可得 $\mu_{\mathcal{A}}(x)\mid\mu_{\mathcal{A},\alpha}(x)$．因此 $\mu_{\mathcal{A},\alpha}(x)=\mu_{\mathcal{A}}(x)$．

\textbf{(2)} 设 $d=\deg\mu_{\mathcal{A},\alpha}(x)$，由 $\mu_{\mathcal{A},\alpha}(x)$ 的定义，结合多项式的带余除法可直接证明
\[
\alpha,\ \mathcal{A}(\alpha),\ \mathcal{A}^2(\alpha),\ \cdots,\ \mathcal{A}^{d-1}(\alpha)
\]
是 $K[\mathcal{A}]\cdot\alpha$ 的一组基．进而 $K[\mathcal{A}]\cdot\alpha$ 是 $\deg\mu_{\mathcal{A},\alpha}(x)$ 维 $\mathcal{A}$-不变子空间．

\textbf{(3)} 由 $\mu_{\mathcal{A}}(x)=\chi_{\mathcal{A}}(x)$，存在 $\alpha\in V$ 使 $V=K[\mathcal{A}]\cdot\alpha$．设 $\mathcal{B}(\alpha)=f(\mathcal{A})(\alpha)$，若 $\mathcal{A}\mathcal{B}=\mathcal{B}\mathcal{A}$，则 $\forall\beta=g(\mathcal{A})(\alpha)\in V$，有
\[
\mathcal{B}(\beta)=g(\mathcal{A})\cdot\mathcal{B}(\alpha)
=g(\mathcal{A})\cdot f(\mathcal{A})(\alpha)=f(\mathcal{A})(\beta).
\]
故 $\mathcal{B}=f(\mathcal{A})$．
\end{solution}

\answerproblem{五、}{15}
设 $f\in C[0,1]\cap C^2(0,1)$，$\sup\limits_{x\in(0,1)}f''(x)=2$．证明：存在唯一的二次多项式 $P(x)=x^2+bx+c$，使得 $P(0)-f(0)=P(1)-f(1)=0$，且下列结论之一成立：
\begin{enumerate}[label=(\arabic*)]
  \item $f(x)>P(x)\ (\forall x\in(0,1))$；
  \item $f(x)=P(x)\ (\forall x\in[0,1])$．
\end{enumerate}
\begin{solution}
\textbf{(i)} 对于 $P(x)=x^2+bx+c$，易见 $P(0)-f(0)=P(1)-f(1)=0$ 成立的充分必要条件是 $c=f(0)$，$b=f(1)-1-f(0)$．因此，满足 $P(0)-f(0)=P(1)-f(1)=0$ 的二次多项式 $P(x)=x^2+bx+c$ 存在唯一．

我们要证明该多项式 $P$ 满足 (1) 或 (2)．

记 $F(x)=P(x)-f(x)$（$x\in[0,1]$），由题设，$F\in C[0,1]\cap C^2(0,1)$，且
\[
F''(x)\geqslant 0,\qquad \forall x\in(0,1).
\]
即 $F$ 是凸函数，因此，
\[
F(x)\leqslant \max\{F(0),F(1)\}=0,\qquad \forall x\in[0,1].
\]

\textbf{(ii)} 我们给出以下两种写法．它们本质上是一样的．

\textbf{写法 I.}\par
若对每个 $x\in(0,1)$，$F(x)=0$，则由连续性，$F$ 在 $[0,1]$ 上恒为零，即情形 (2) 发生．

否则，存在一个 $\xi\in(0,1)$ 使得 $F(\xi)<0$．从而，对于 $x\in(0,\xi)$，有
\[
F(x)=F\Bigl(\frac{x}{\xi}\,\xi\Bigr)
\leqslant\frac{x}{\xi}F(\xi)+\Bigl(1-\frac{x}{\xi}\Bigr)F(0)<0.
\]
而对于 $x\in(\xi,1)$，记 $\alpha=\dfrac{1-x}{1-\xi}$，有
\[
F(x)=F\bigl(\alpha\xi+(1-\alpha)\bigr)
\leqslant\alpha F(\xi)+(1-\alpha)F(1)<0.
\]
因此，
\[
F(x)<0,\qquad \forall x\in(0,1).
\]
即情形 (1) 发生．

\textbf{写法 II.}\par
若对每个 $x\in(0,1)$，$F(x)<0$，则情形 (1) 发生．

否则，存在一个 $\xi\in(0,1)$ 使得 $F(\xi)=0$．从而，对于 $x\in(0,\xi)$，记 $\alpha=\dfrac{1-\xi}{1-x}$，有
\[
F(x)=\frac{1}{\alpha}\bigl(\alpha F(x)+(1-\alpha)F(1)\bigr)
\geqslant\frac{1}{\alpha}F\bigl(\alpha x+(1-\alpha)\cdot 1\bigr)
=\frac{1}{\alpha}F(\xi)=0.
\]
而对于 $x\in(\xi,1)$，有
\[
F(x)=\frac{x}{\xi}\Bigl(\frac{\xi}{x}F(x)+\frac{x-\xi}{x}F(0)\Bigr)
\geqslant\frac{x}{\xi}F(\xi)=0.
\]
因此，
\[
F(x)\geqslant 0,\qquad \forall x\in(0,1).
\]
结合 $F\leqslant 0$ 以及连续性得到 $F$ 在 $[0,1]$ 上恒为零．即情形 (2) 发生．$\square$

\textbf{注 1.} 证明中直接使用闭区间 $[a,b]$ 上凸函数 $g$ 的如下一些众所周知的基本性质，不扣分：

(i) 凸函数的最大值在边界取到．

(ii) 可导凸函数的导数单增．

(iii) 对于 $a\leqslant x<y<z\leqslant b$，
\[
\frac{g(y)-g(x)}{y-x}\leqslant\frac{g(z)-g(x)}{z-x}\leqslant\frac{g(z)-g(y)}{z-y}.
\]

\textbf{注 2.} 本题证明的写法可以有非常多的变化．比如，若在 $x_0\in(0,1)$ 处，$F(x_0)=0$，

利用 (ii)，则 $F$ 在 $x_0$ 处取得极大值，从而 $F'(x_0)=0$．因此，$F'$ 在 $(0,x_0)$ 内小于等于零，在 $(x_0,1)$ 内大于等于零，于是 $x_0$ 又是 $F$ 的最小值点，从而 $F\equiv 0$．

利用 (iii)，则对任何 $0\leqslant x<x_0<y\leqslant 1$，有
\begin{align*}
\frac{F(x)-F(x_0)}{x-x_0}&\leqslant\frac{F(1)-F(x_0)}{1-x_0}=0,\\
\frac{F(y)-F(x_0)}{y-x_0}&\leqslant\frac{F(x_0)-F(0)}{x_0-0}=0.
\end{align*}
从而结合 $F\leqslant 0$ 可得 $F(x)=F(y)=0$．
\end{solution}

\answerproblem{六、}{20}
设 $\mathcal{A}$ 为正整数集 $\mathbb{Z}_+$ 的子集，$\lim\limits_{n\to+\infty}\frac{S_n(0,1)}{n}$ 存在，记为 $p$，其中对于 $0\leqslant a\leqslant b\leqslant 1$，$S_n(a,b)$ 表示集合 $\{\,k\in\mathcal{A}\mid an\leqslant k\leqslant bn\,\}$ 的元素个数．
\begin{enumerate}[label=(\arabic*)]
  \item 若 $0<a<b<1$，证明极限 $\lim\limits_{n\to+\infty}\frac{S_n(a,b)}{n}$ 存在并求其值．
  \item 若 $f$ 在 $[0,1]$ 上 \textbf{Riemann} 可积，证明：
  \[
    \lim_{n\to+\infty}\frac{1}{n}\sum_{k\leqslant n,\ k\in\mathcal{A}}f\Bigl(\frac{k}{n}\Bigr)=p\int_0^1 f(x)\dd x.
  \]
\end{enumerate}
\begin{solution}
\textbf{(1)} 考虑 $n>\dfrac{2}{a}$．用 $[x]$ 表示实数 $x$ 的整数部分，则易见
\begin{align*}
\{k\in\mathcal{A}\mid an\leqslant k\leqslant bn\}
&\subseteq\{k\in\mathcal{A}\mid [an]-1<k\leqslant [bn]+1\}\\
&=\{k\in\mathcal{A}\mid 1\leqslant k\leqslant [bn]+1\}\setminus\{k\in\mathcal{A}\mid 1\leqslant k\leqslant [an]-1\},
\end{align*}
因此，
\[
S_n(a,b)\leqslant S_{[bn]+1}(0,1)-S_{[an]-1}(0,1),
\]
类似地，有
\[
S_n(a,b)\geqslant S_{[bn]}(0,1)-S_{[an]+1}(0,1),
\]
因此，由题设条件可得
\[
\lim_{n\to+\infty}\frac{S_n(a,b)}{n}=(b-a)p.
\]

\textbf{(2)} 由 (1) 的结论，对于 $0<a<b<1$ 以及 $f=\chi_{[a,b]}$，本小题的结论成立：
\[
\lim_{n\to+\infty}\frac{1}{n}\sum_{k\leqslant n,\ k\in\mathcal{A}}f\Bigl(\frac{k}{n}\Bigr)=p\int_0^1 f(x)\dd x.
\]
事实上，仿 (1) 可证对于 $0\leqslant a\leqslant b\leqslant 1$ 以及 $f=\chi_{[a,b]},\chi_{(a,b]},\chi_{[a,b)},\chi_{(a,b)}$，结论成立（$[a,a]$ 理解为 $\{a\}$）．进而当 $f$ 是 $[0,1]$ 上分段常值函数时，本小题的结论成立．

一般地，设 $f$ 在 $[0,1]$ 上 Riemann 可积，设 $M$ 为 $|f|$ 的一个上界．任取 $m\geqslant 1$，将 $[0,1]$ 作 $m$ 等分，记该等分划分为 $P_m$，定义
\[
f_m(x):=\begin{cases}
\sup\limits_{t\in\left[\frac{k}{m},\frac{k+1}{m}\right]}f(t), & x\in\left(\frac{k}{m},\frac{k+1}{m}\right),\ 0\leqslant k\leqslant m-1;\\[2mm]
M, & \text{其他}.
\end{cases}
\]
\[
g_m(x):=\begin{cases}
\inf\limits_{t\in\left[\frac{k}{m},\frac{k+1}{m}\right]}f(t), & x\in\left(\frac{k}{m},\frac{k+1}{m}\right),\ 0\leqslant k\leqslant m-1;\\[2mm]
-M, & \text{其他}.
\end{cases}
\]
则
\[
g_m(x)\leqslant f(x)\leqslant f_m(x),\qquad \forall x\in[0,1],
\]
\[
\int_0^1 g_m(x)\dd x=L(f;P_m),\qquad \int_0^1 f_m(x)\dd x=U(f;P_m),
\]
其中 $U(f;P_m)$，$L(f;P_m)$ 分别表示 $f$ 对应于划分 $P_m$ 的 Darboux 上和与下和．这样，
\[
\frac{1}{n}\sum_{k\leqslant n,\ k\in\mathcal{A}}g_m\Bigl(\frac{k}{n}\Bigr)
\leqslant\frac{1}{n}\sum_{k\leqslant n,\ k\in\mathcal{A}}f\Bigl(\frac{k}{n}\Bigr)
\leqslant\frac{1}{n}\sum_{k\leqslant n,\ k\in\mathcal{A}}f_m\Bigl(\frac{k}{n}\Bigr),
\]
并令 $n\to+\infty$ 得到
\[
pL(f;P_m)\leqslant\varliminf_{n\to+\infty}\frac{1}{n}\sum_{k\leqslant n,\ k\in\mathcal{A}}f\Bigl(\frac{k}{n}\Bigr)
\leqslant\varlimsup_{n\to+\infty}\frac{1}{n}\sum_{k\leqslant n,\ k\in\mathcal{A}}f\Bigl(\frac{k}{n}\Bigr)
\leqslant pU(f;P_m).
\]
再令 $m\to+\infty$ 即得
\[
\lim_{n\to+\infty}\frac{1}{n}\sum_{k\leqslant n,\ k\in\mathcal{A}}f\Bigl(\frac{k}{n}\Bigr)=p\int_0^1 f(x)\dd x.
\]
结论得证．$\square$

\textbf{注.} 若第二小题采用其他证明途径．

对于 $f(x)=x^k$ 或 $f(x)=\cos 2k\pi x$，或 $f(x)=\sin 2k\pi x$，证明结论成立，可考虑在第一小题 $5$ 分基础上再给 $5$ 分．

若进一步证明 $f$ 为连续函数（连续周期函数），或 $f$ 为分段线性的连续函数时，结论成立，则可以共给到 $14$ 分．
\end{solution}
